(1) 记 \(a=n_1\Phi_0-\Phi_{e1}\)、\(b=n_2\Phi_0-\Phi_{e2}\)，磁通约束等价于 \(\begin{pmatrix}L&M\\M&L\end{pmatrix}(I_1,I_2)^T=(a,b)^T\)。逆矩阵存在因为 \(L^2-M^2>0\)，求得 \(I_1=(La-Mb)/(L^2-M^2)\)、\(I_2=(Lb-Ma)/(L^2-M^2)\)。把两个环误当作无互感会丢掉一个回路改变另一回路持久电流的效应。

(2) 电感能量 \(U=\tfrac12 L(I_1^2+I_2^2)+MI_1I_2\)。变到 \(I_\pm=(I_1\pm I_2)/\sqrt2\)，得 \(U=\tfrac12(L+M)I_+^2+\tfrac12(L-M)I_-^2\)。因 \(|M|<L\)，两模系数均正，非零环流能量严格正；若 \(|M|=L\)，电感矩阵奇异，意味着理想独立回路模型退化，不能直接反演。

(3) 当 \(\Phi_{e1}=\Phi_{e2}=\Phi_e,\ n_1=n_2=n\)，有 \(a=b=n\Phi_0-\Phi_e\)，从而 \(I_1=I_2=(n\Phi_0-\Phi_e)/(L+M)\)。能量为 \(U=(n\Phi_0-\Phi_e)^2/(L+M)\)。在这个对称子空间，磁通偏离每个量子数 \(n\Phi_0\) 时产生回复的超导电流；但非零互感改变回复系数。

(4) 固定整数 \(n_1,n_2\) 时，\(\partial I_1/\partial\Phi_{e1}=-L/(L^2-M^2)\)、\(\partial I_1/\partial\Phi_{e2}=M/(L^2-M^2)\)，其余由对称性获得。允许相位滑移后才可在不同整数态间重新选择较低自由能，可能出现持久电流的阶跃与滞后；不能在同一绝热分支中把 \(n_i\) 当作连续实数来抵消任意外磁通。磁能表达式是固定流态的部分能量，外场线圈作功另计。
